\documentclass[10pt]{article}
\usepackage[margin=1in]{geometry}
\usepackage{amsmath,amssymb}
\usepackage{graphicx}
\usepackage{booktabs}
\usepackage{hyperref}
\usepackage{natbib}

\title{Supplementary Materials: Transition-Field Geometry Across Heterogeneous Dynamical Systems}
\author{Mark Rowe Traver}
\date{}

\begin{document}
\maketitle

\section{Null Model Details}
\label{sec:null_details}

\subsection{N1: IID Gaussian}

For each feature $j \in \{1, \dots, 17\}$, we sample $X_{\text{null},ij} \sim \mathcal{N}(\mu_j, \sigma_j^2)$ independently, where $\mu_j$ and $\sigma_j^2$ are the empirical mean and variance of feature $j$ across all 211 systems. This destroys all inter-feature correlations while preserving marginal distributions.

\subsection{N2: Covariance-Preserving Gaussian}

We sample $\mathbf{X}_{\text{null}} \sim \mathcal{N}(\boldsymbol{\mu}, \boldsymbol{\Sigma})$ where $\boldsymbol{\mu} \in \mathbb{R}^{17}$ is the empirical mean and $\boldsymbol{\Sigma} \in \mathbb{R}^{17 \times 17}$ is the empirical covariance matrix of the 211 systems. This preserves all second-order statistics.

\subsection{N3: Spectrum-Preserving Randomization}

Let $\Phi = \mathbf{U} \mathbf{S} \mathbf{V}^\top$ be the SVD of the centered $\Phi$-coordinates. We generate random orthogonal matrices $\mathbf{U}'$ and $\mathbf{V}'$ via QR decomposition of random Gaussian matrices, and construct $\Phi_{\text{null}} = \mathbf{U}' \mathbf{S} \mathbf{V}'^\top$. This preserves the singular value spectrum while randomizing the orientation.

\subsection{N4: Copula-Preserving Shuffle}

\begin{enumerate}
\item Compute empirical CDF of each feature: $F_j(x) = \frac{1}{N} \sum_{i=1}^N \mathbf{1}[x_{ij} \leq x]$.
\item Transform to uniform marginals: $U_{ij} = F_j(x_{ij})$.
\item Compute the Gaussian copula: sample $\mathbf{Z} \sim \mathcal{N}(\mathbf{0}, \mathbf{R})$ where $\mathbf{R}$ is the rank correlation matrix.
\item Transform to uniform: $\tilde{U}_{ij} = \Phi(Z_{ij})$ (standard normal CDF).
\item Inverse-transform to marginals: $X_{\text{null},ij} = F_j^{-1}(\tilde{U}_{ij})$.
\end{enumerate}

This preserves marginal distributions and rank correlations while destroying the specific pairing of values.

\subsection{N5: Maximum Entropy}

Equivalent to N2 under second-moment constraints. The maximum entropy distribution with given mean and covariance is the Gaussian, so N5 = N2. This explicitly tests whether matching first and second moments is sufficient.

\subsection{N6: Adversarial Low-Rank}

We construct a synthetic manifold in $\mathbb{R}^4$ with matched participation ratio:
\begin{enumerate}
\item Compute the target PR from the real $\Phi$-coordinates.
\item Construct a diagonal covariance matrix with eigenvalues $\lambda_k \propto e^{-k/\alpha}$ where $\alpha$ is chosen to match the target PR.
\item Generate random orthogonal eigenvectors via QR decomposition.
\item Sample from the resulting Gaussian mixture.
\end{enumerate}

\subsection{N7: Random Embedding}

We apply a random projection matrix $\mathbf{R} \in \mathbb{R}^{17 \times 17}$ (with orthonormal columns) to the feature matrix: $\mathbf{X}_{\text{null}} = \mathbf{X} \mathbf{R}$. This tests whether the emergence coordinate transformation is sensitive to the specific feature ordering.

\section{Manifold Metrics}
\label{sec:metrics}

\subsection{Participation Ratio}

For a point cloud $\Phi \in \mathbb{R}^{N \times d}$ with covariance matrix $\mathbf{C} = \frac{1}{N} \sum_{i=1}^N (\phi_i - \bar{\phi})(\phi_i - \bar{\phi})^\top$ and eigenvalues $\{\lambda_k\}$, the participation ratio is:
\[
\text{PR} = \frac{(\sum_k \lambda_k)^2}{\sum_k \lambda_k^2}
\]
PR ranges from 1 (all variance in one dimension) to $d$ (uniform variance across all dimensions).

\subsection{kNN Flow Prediction}

For each system $i$, we predict its flow magnitude as the average flow magnitude of its $k$-nearest neighbors in $\Phi$-space:
\[
\hat{v}_i = \frac{1}{k} \sum_{j \in \text{NN}_k(i)} v_j
\]
The kNN flow prediction score is the Pearson correlation between $\{\hat{v}_i\}$ and $\{v_i\}$ across all systems.

\subsection{Smoothness}

The smoothness metric measures the local regularity of the manifold by computing the mean magnitude of the gradient of the log-density along the flow direction:
\[
S = \frac{1}{N} \sum_{i=1}^N \left\| \frac{\nabla \log \rho(\phi_i)}{\|\nabla \log \rho(\phi_i)\|} \cdot \hat{\mathbf{v}}_i \right\|
\]
where $\rho(\phi)$ is the kernel density estimate and $\hat{\mathbf{v}}_i$ is the unit flow vector.

\section{Bootstrap Confidence Intervals}
\label{sec:bootstrap}

All confidence intervals are computed via nonparametric bootstrap with $B = 100$ resamples. For a metric $m$, we compute:
\[
\hat{m}_b = m(\Phi_{\text{bootstrap}}), \quad b = 1, \dots, B
\]
where $\Phi_{\text{bootstrap}}$ is obtained by sampling $N$ systems with replacement. The 95\% confidence interval is:
\[
[\hat{m}_{(\alpha/2)}, \hat{m}_{(1-\alpha/2)}]
\]
where $\hat{m}_{(p)}$ is the $p$-th empirical quantile.

\section{Feature Definitions}
\label{sec:features}

\begin{table}[h]
\centering
\caption{The 17 raw features used in the transition-field analysis.}
\begin{tabular}{ll}
\toprule
Feature & Definition \\
\midrule
pc1 & First principal component variance \\
pc2 & Second principal component variance \\
effective_rank & $\frac{(\sum \sigma_i)^2}{\sum \sigma_i^2}$ \\
$\tau_{m1}$--$\tau_{m4}$ & Moment-based timescales \\
temporal_corr & Temporal autocorrelation \\
phase_corr & Phase-space correlation \\
pc1\_ratio & PC1 variance ratio \\
replay\_displacement & Replay displacement metric \\
abl\_full\_pc1 & Full ablation sensitivity \\
abl\_no\_m1\_pc1 & Ablation without $m_1$ \\
abl\_no\_m2\_pc1 & Ablation without $m_2$ \\
abl\_no\_m3\_pc1 & Ablation without $m_3$ \\
abl\_no\_m4\_pc1 & Ablation without $m_4$ \\
m2\_contribution & $m_2$ contribution fraction \\
\bottomrule
\end{tabular}
\end{table}

\section{System Families}
\label{sec:families}

\begin{table}[h]
\centering
\caption{The 13 dynamical families and their parameter ranges.}
\begin{tabular}{lcc}
\toprule
Family & Systems & Parameter Range \\
\midrule
Logistic & 16 & $r \in [2.5, 4.0]$ \\
Lorenz & 16 & $\sigma \in [8, 16]$ \\
R\"ossler & 16 & $a \in [0.1, 0.5]$ \\
H\'enon & 16 & $a \in [0.5, 1.5]$ \\
H\'enon--Heiles & 16 & $E \in [0.05, 0.20]$ \\
H\'enon Map & 16 & $a \in [0.5, 1.5]$ \\
Ikeda & 16 & $u \in [0.3, 1.0]$ \\
Duffing & 16 & $\alpha \in [-0.5, 0.5]$ \\
Van der Pol & 16 & $\mu \in [0.5, 5.0]$ \\
Chua & 16 & $m_0 \in [-1.2, -0.6]$ \\
Brusselator & 16 & $A \in [1, 3]$ \\
Stuart--Landau & 16 & $\lambda \in [-0.5, 0.5]$ \\
FitzHugh--Nagumo & 16 & $a \in [0.5, 1.5]$ \\
\bottomrule
\end{tabular}
\end{table}

\section{Reproducibility}
\label{sec:reproducibility}

All analyses were performed with deterministic random seed 42. The complete pipeline is:

\begin{verbatim}
python T031_NULL_GEOMETRY_DISENTANGLEMENT.py
\end{verbatim}

Runtime: approximately 7 seconds on standard hardware (Python 3.14, NumPy, SciPy, scikit-learn).

Environment: see \texttt{reproducibility/environment.yml} for exact package versions.

\end{document}
